\section{The weakness taxonomy W1--W14}\label{sec:taxonomy}

Table~\ref{tab:taxonomy} defines the fourteen weakness classes used through the
rest of the paper. Anchors: OWASP Top 10 for LLM Applications (2026 agentic
edition) \cite{owasp2026} and MITRE ATLAS techniques \cite{atlas}. The \emph{lens}
column carries the paper's central discipline: \textbf{O} = offensive class (we
wield it against targets), \textbf{D} = defensive exposure (our own stack is the
victim), \textbf{B} = dual-use (both).

\begin{longtable}{L{0.8cm} L{3.0cm} L{4.2cm} L{3.1cm} L{0.8cm}}
\caption{The fourteen weakness classes. OWASP codes reference the 2026 agentic
edition; ATLAS codes the 2026.10 matrix.}\label{tab:taxonomy}\\
\toprule
ID & Class & Definition & Anchors & Lens \\
\midrule
\endfirsthead
\toprule
ID & Class & Definition & Anchors & Lens \\
\midrule
\endhead
\bottomrule
\endfoot
W1 & Injection vs defended stacks & Indirect prompt injection against agents that
already deploy injection defenses; the residual rate that survives tool firewalls
and spotlighting. & LLM01; AML.T0051 & D \\
W2 & Adaptive jailbreak & Iterative, response-aware jailbreak of the agent stack
itself (PAIR/TAP-style escalation), including against our guard chain. &
LLM01; AML.T0048 & B \\
W3 & Swarm cascade & Compromise propagating across cooperating agents (payloads in
inter-agent messages, infected shared context). & AML.T0051 agentic; MAST & O \\
W4 & Tool/MCP supply chain & Poisoned tool descriptions, malicious MCP servers,
fuzzable tool registries; the ToolFuzz surface. & LLM06 (supply chain);
AML tool-exploitation family & D \\
W5 & Memory poisoning & Adversarial writes to long-term memory/experience stores
that persist across runs and reshape future behavior. & Agentic memory class
(AgentPoison family) & D \\
W6 & Hidden-context exfiltration & Data leaving through render/egress channels
(reference-style Markdown, auto-fetched images) driven by retrieved instructions
--- the EchoLeak shape. & LLM02/hidden-context exposure; CWE-200 & B \\
W7 & Retry-inflated measurement & Measurement validity: retry-until-success
harnesses reporting session rates as per-attempt capability. & Evaluation-method
class (this paper) & B \\
W8 & Judge miscalibration & LLM-as-judge bias inflating or deflating ASR and
safety scores. & Evaluation-method class \cite{judgesurvey} & B \\
W9 & Single-seed reliability & Reporting pass@1 as capability when pass$^k$ across
independent runs collapses (tau-bench pass$^8<25\%$ class). &
Evaluation-method class \cite{taubench} & B \\
W10 & Cost/turn amplification & Agentic loops turned into budget-exhaustion DoS;
cost amplification as an attack primitive. & Availability class & O \\
W11 & Misuse compliance & Agents complying with overtly harmful tasks when asked
politely (AgentHarm/SafeArena rates). & LLM01; AML.T0048 & O \\
W12 & Output-chain second order & Our outputs becoming another system's inputs
with unverified trust; downstream consumers as exfil/enrichment channels. &
LLM06?---inter-system; CWE-200 & B \\
W13 & Multilingual erosion & Guardrail strength decaying outside English;
multilingual jailbreaks and safety drift. & LLM01 & O \\
W14 & Supply-chain/model drift & Backend/model/pin drift changing behavior under
the same code; silent capability shifts. & LLM06?---supply chain & D \\
\end{longtable}

\subsection{Reading the 2026 OWASP reshuffle}\label{sec:owasp2026}

Three moves in the 2026 edition matter more than the rest. First, \emph{excessive
agency} at position~3 ratifies what the incident record already showed: the harm
comes from \emph{what the system is allowed to do}, not only from what it says.
Second, \emph{system prompt leakage} becoming \emph{hidden context exposure}
widens the class from ``the prompt leaked'' to ``retrieved context did something
the operator never wrote''---which is EchoLeak's exact structure: our own
retrieval became the attack delivery channel. Third, the list is now written for
systems that \emph{act}: tool effects, cascading actions, and cross-system trust
appear as first-class categories rather than footnotes.

Our own incident log contains a miniature of this reshuffle. In an internal
exercise, a red-team finding was stored in two halves: a screenshot of a vulnerable
page went into one store, and the rendered proof text went through a rendering
path that treated it as content. Neither half was sensitive alone; the recombined
pair was a working exfiltration story. That is W6 in one paragraph---and it is
why the D-lens exists: the framework's own storage and rendering pipeline recreated
the vulnerability class it was built to hunt.
